\documentclass[../../../include/open-logic-section]{subfiles}

\begin{document}

\olfileid{his}{set}{pathology}
\olsection{વિચિત્ર રચનાઓ}

પરંતુ \emph{લક્ષ}ની વ્યાખ્યાએ કેટલીક ખૂબ
“વિચિત્ર” રચનાઓને પણ સ્થાન આપ્યું.

1830ના દાયકાની આસપાસ બોલ્ઝાનોએ એવું
વિધેય શોધ્યું જે \emph{દરેક જગ્યાએ સતત}
પણ \emph{ક્યાંય વિકલનીય નહીં} હતું.
(દુર્ભાગ્યે બોલ્ઝાનોએ આ શોધ પ્રકાશિત
કરી નહોતી; 1872માં વાયરસ્ટ્રાસે સ્વતંત્ર
રીતે એ જ વિચાર શોધ્યો ત્યારે ગણિતશાસ્ત્રીઓને
તે પહેલી વાર જાણ થયો.)\footnote{આ ઇતિહાસ
\href{http://en.wikipedia.org/wiki/Weierstrass_function}{વાયરસ્ટ્રાસ
વિધેય} અંગેના વિકિપીડિયા લેખની અત્યંત
વિગતવાર પાદટીપોમાં નોંધાયેલો છે.}
આ, ઓછામાં ઓછું કહીએ તો, ખૂબ આશ્ચર્યજનક
હતું. $|x|$ જેવું વિધેય શોધવું સહેલું છે
જે દરેક જગ્યાએ સતત હોય પણ કોઈ એક
બિંદુએ વિકલનીય ન હોય. પરંતુ દરેક
જગ્યાએ સતત અને \emph{ક્યાંય}
વિકલનીય ન હોય એવું વિધેય સાવ
જુદી બાબત છે. જરા વિચારો કે
આવું વિધેય કેવી રીતે દોરશો.
સતતતા માટે કાગળ પરથી પેન
ઉઠાવ્યા વિના રેખા દોરી શકાય
એવું હોવું જોઈએ; પરંતુ ક્યાંય
વિકલનીય ન હોય તે માટે પેનની
દિશા સતત એકાએક બદલવી પડે.

પછી આવી બીજી “વિચિત્ર” બાબતો સામે આવી.
5 જાન્યુઆરી 1874ના રોજ કૅન્ટૉરે
ડેડેકિન્ડને પત્ર લખીને આ પ્રશ્ન મૂક્યો:
\begin{quote}
શું કોઈ સપાટીને (દાખલા તરીકે, તેની સીમા
સહિતના ચોરસને) કોઈ રેખા સાથે
(દાખલા તરીકે, અંતબિંદુઓ સહિતની
સીધી રેખા સાથે) એક-એક સંગતતામાં
મૂકી શકાય, જેથી સપાટીના દરેક
બિંદુને રેખાનું એક બિંદુ અને,
ઊલટું, રેખાના દરેક બિંદુને
સપાટીનું એક બિંદુ અનુરૂપ હોય?

હજુ મને આ પ્રશ્નનો ઉત્તર બહુ અઘરો લાગે
છે—જોકે અહીં પણ \emph{ના} કહેવાની
ઇચ્છા એટલી પ્રબળ થાય છે કે તેની
સાબિતી લગભગ અનાવશ્યક ગણવાનું
મન થાય. [\citealt{Gouvea2011}માં
ઉદ્ધૃત]
\end{quote}
પરંતુ 1877માં કૅન્ટૉરે પોતે ખોટા
હોવાની સાબિતી આપી. હકીકતમાં
રેખા અને ચોરસમાં બિંદુઓની
સંખ્યા બરાબર સમાન છે.
29 જૂન 1877ના રોજ તેમણે
ડેડેકિન્ડને લખ્યું:
“\emph{je le vois, mais je ne le crois pas}”;
એટલે, “હું તે જોઈ શકું છું,
પણ તેનો વિશ્વાસ નથી થતો.”
ગણિતના “સ્વીકૃત ઇતિહાસ”માં
આ વાક્ય ઘણી વાર એ દર્શાવતું
માનવામાં આવે છે કે તે સમયના
ગણિતશાસ્ત્રીઓને આ નવાં
પરિણામો \emph{શાબ્દિક અર્થમાં
અવિશ્વસનીય} લાગતાં હતાં.
(પત્રવ્યવહાર \citet{Gouvea2011}માં
આપેલો છે; આપણે
\olref[his][set][mythology]{sec}માં
તેની તરફ ફરી વળીશું.
કૅન્ટૉરની સાબિતીનો
ખાકો \olref[his][set][cantorplane]{sec}માં
આપ્યો છે.)

કૅન્ટૉરના પરિણામથી પ્રેરાઈ પીઆનોએ
વિચાર્યું કે રેખાને સમતલ પર
\emph{સાતત્યપૂર્વક} ચિતરી શકાય કે કેમ.
\sourcecorrection{OLMD-168}{મૂળમાં
\texttt{smoothly} શબ્દ છે. અહીં
ગાણિતિક અર્થમાં વિકલનીય અથવા
મસૃણ ચિત્રણ શક્ય નથી;
પીઆનોનું અવકાશપૂરક ચિત્રણ
સતત છે. તેથી “સાતત્યપૂર્વક”
લખ્યું છે.}
આ \emph{અવકાશ ભરતી વક્રરેખા}
હશે. \citeyear{Peano1890}માં
પીઆનોએ આવી જ વક્રરેખા
રચી. આ ખરેખર સહજ
સમજને પડકારતું છે:
યુક્લિડે રેખાને
“પહોળાઈ વિનાની
લંબાઈ” કહેલી હતી
(પુસ્તક ૧, વ્યાખ્યા ૨),
પરંતુ પીઆનોએ બતાવ્યું
કે રેખાને યોગ્ય રીતે
વાળી તેની લંબાઈથી
પહોળાઈ ભરી શકાય.
\citeyear{Hilbert1891}માં
હિલ્બર્ટે વધુ સહજ
લાગે એવી અવકાશપૂરક
વક્રરેખા અને તેને
દર્શાવતાં ચિત્રો આપ્યાં.
આ વક્રરેખા ક્રમશઃ
રચાય છે; તેના
પ્રથમ છ તબક્કા
અહીં છે:
\begin{center}
\begin{tikzpicture}
\begin{scope}[xshift=20pt, yshift=20pt]
	\draw [oldiagcolorC, l-system={Hilbert curve, step=40pt, angle=90, axiom=L, order=1}] lindenmayer system;
\end{scope}
\begin{scope}[xshift=110pt, yshift=10pt]
	\draw [oldiagcolorC, l-system={Hilbert curve, step=20pt, angle=90, axiom=L, order=2}] lindenmayer system;
\end{scope}
\begin{scope}[xshift=205pt, yshift=5pt]
	\draw [oldiagcolorC, l-system={Hilbert curve, step=10pt, angle=90, axiom=L, order=3}] lindenmayer system;
\end{scope}
\begin{scope}[xshift=2.5pt, yshift=-97.5pt]
	\draw [oldiagcolorC, l-system={Hilbert curve, step=5pt, angle=90, axiom=L, order=4}] lindenmayer system;
\end{scope}
\begin{scope}[xshift=101.25pt, yshift=-98.75pt]
	\draw [oldiagcolorC, l-system={Hilbert curve, step=2.5pt, angle=90, axiom=L, order=5}] lindenmayer system;
\end{scope}
\begin{scope}[xshift=200.625pt, yshift=-99.375pt]
	\draw [oldiagcolorC, l-system={Hilbert curve, step=1.25pt, angle=90, axiom=L, order=6}] lindenmayer system;
\end{scope}
\draw[gray] (0pt,0pt) rectangle (80pt, 80pt);
\draw[gray] (100pt,0pt) rectangle (180pt, 80pt);
\draw[gray] (200pt,0pt) rectangle (280pt, 80pt);
\draw[gray] (0pt,-100pt) rectangle (80pt, -20pt);
\draw[gray] (100pt,-100pt) rectangle (180pt, -20pt);
\draw[gray] (200pt,-100pt) rectangle (280pt, -20pt);
\end{tikzpicture}  
\end{center}
લક્ષ લેતાં—અને હવે લક્ષની ચુસ્ત
વ્યાખ્યા મળી ગઈ છે—આખો
ચોરસ ગાઢ !!{colorC} રંગથી
ભરાઈ જાય છે. વળી,
હિલ્બર્ટની વક્રરેખા
દરેક જગ્યાએ સતત પણ
ક્યાંય વિકલનીય નથી;
સહજ રીતે કહીએ તો,
અનંત લક્ષમાં વિધેય
દરેક ક્ષણે દિશા
એકાએક બદલે છે.
(હિલ્બર્ટની રચનાનો
વધુ વિગતવાર ખાકો
\olref[his][set][hilbertcurve]{sec}માં
આપીશું.)

ભલે સારું કે ખરાબ, આ “વિચિત્ર”
ભૌમિતિક રચનાઓને ભૌમિતિક
સહજજ્ઞાન પર આધાર રાખવા
સામે શંકાનું કારણ માનવામાં
આવી. તે ગણિત કરવાની
નવી રીત માટે લગભગ
\emph{પ્રચાર} બની,
જે આગળ ગણસિદ્ધાંતમાં
પરિણમી. વિષય વિશેની
પછીથી રચાયેલી
દંતકથાઓમાં વારંવાર
કહેવામાં આવતું કે
આ પરિણામો સંપૂર્ણ
ચુસ્ત પણ હતાં અને
સંપૂર્ણ આશ્ચર્યજનક
પણ. તેથી તેમણે
બેવડો હેતુ સાધ્યો:
ભૌમિતિક સહજજ્ઞાન
પર નિર્ભર ન રહેવાની
ચેતવણી આપી અને
નવી વિચારરીતોની
ફળદ્રુપતા બતાવી.

\end{document}
